\section{合规、审计与工程落地}
\subsection{Logging 与可追溯性}
每个 policy update 应该伴随日志：seed、hyper parameters、KL divergence、critic loss。结论应该链接到 `audit note` 以便 regulator 文件化 assessments。

\begin{table}[H]
\centering
\small
\begin{tabular}{p{0.3\textwidth}p{0.6\textwidth}}
\toprule
Artifact & 内容 \\
\midrule
Prompt registry & 记录 prompt template + temperature \\
Override log & 记录 clinician override + rationale \\
Evidence patch bank & 存储 attention map 与 highlight frames \\
Release ticket & 包含 rollout plan + rollback criteria \\
\bottomrule
\end{tabular}
\caption{Policy gradient release 时的 audit artifacts}
\end{table}

\subsection{工程交付 checklist}
每次 release 前必须 complete：
\begin{enumerate}
    \item 冻结 prompt + hyper parameters；
    \item 复核 KL divergence < threshold；
    \item 归档 reward curve + entropy log；
    \item 记录 fail-safe policy + rollback plan。
\end{enumerate}

\begin{importantbox}{Release checklist 的价值}
\begin{itemize}
    \item 确保每次 update 都有可追溯的决策路径；
    \item 加速 regulator review 过程；
    \item 形成持续改进的 knowledge base。
\end{itemize}
\end{importantbox}

\subsection{本章小结}
合规与 audit 是 policy gradient 落地的最后一公里：日志、artifact、checklist 构成了 regulator 可接受的 evidence trail。

